%==============================================================
% BAB II  TINJAUAN PUSTAKA (OPSIONAL)
%==============================================================

\chapter{TINJAUAN PUSTAKA (OPSIONAL)}

Pustaka yang digunakan dalam bab ini ialah acuan primer, diutamakan artikel
jurnal dan paten yang relevan dengan bidang yang diteliti, terkini, dan asli
(\textit{state of the art}).
Diktat dan buku ajar tidak termasuk acuan primer.
Tinjauan pustaka memuat telaah singkat, jelas, dan sistematis tentang kerangka
teoretis, kerangka pikir, temuan, postulat-postulat, prinsip, asumsi, dan
hasil-hasil penelitian yang relevan yang melandasi masalah penelitian atau
gagasan guna menggali pemahaman mengenai masalah penelitian dan pemecahan
masalahnya.
Oleh karena itu, dari tinjauan pustaka harus dapat diturunkan kerangka pikir,
hipotesis penelitian, dan metode penelitian.

%--------------------------------------------------------------
% CONTOH TABEL KEDUA (lanjutan sequential dari Tabel 1)
% Tabel ini menunjukkan cara tabel dengan baris subgrup dan
% catatan kaki superskrip.
%--------------------------------------------------------------

\begin{table}[htbp]
  \caption{Tingkat kekerasan buah pisang raja pada suhu simpan yang berbeda
           dan pemberian putresina}
  \label{tab:pisang_raja}
  \centering
  \begin{tabular}{lcccc}
    \toprule
\multirow{3}{*}{Perlakuan}
  & \multicolumn{4}{c}{Kekerasan buah dan kandungan gula pada hari ke-} \\
\cmidrule(lr){2-5}
  & 0 &  & 7 & 14 \\
\midrule
  & \multicolumn{4}{c}{Kekerasan buah (mm 50 g$^{-1}$ detik$^{-1}$)$^{\mathrm{a}}$} \\
    \multicolumn{5}{l}{\textit{Suhu Simpan}} \\
    \quad 15 ºC              & 9.20a  &  & 13.40a & 11.83a \\
    \quad 28 ºC              & 10.64a &  & 11.22a & 80.43b \\
    \multicolumn{5}{l}{\textit{Putresina}} \\
    \quad Dengan putresina   & 12.07a &  & 13.23a & 11.19a \\
    \quad Tanpa putresina    & 10.76a &  & 14.41a & 41.12b \\
    \bottomrule
  \end{tabular}
  \begin{minipage}{\linewidth}
    \vspace{2pt}
    \footnotesize
    $^{\mathrm{a}}$Angka-angka pada kolom yang sama yang diikuti oleh huruf yang
    sama tidak berbeda nyata pada taraf uji 5\% (uji selang berganda Duncan).
  \end{minipage}
\end{table}


%--------------------------------------------------------------
\section{Contoh Subbab}
%--------------------------------------------------------------

% --------------------------------------------------------------
\subsection{Contoh Sub-subbab}
% --------------------------------------------------------------

Berikut adalah contoh sub-subbab. Pada sub-subbab ini posisi paragraf lebih menjorok dari paragraf di subbab

%--------------------------------------------------------------
\section{Contoh Subbab2}
%--------------------------------------------------------------

Contoh tulisan